\section{Introduction}
\label{sec:introduction}
Understanding biological systems requires reasoning across relationships and scales: from molecular interactions and cellular organization to neural connectivity and coordinated physiological processes. Visual analysis makes these relationships available for inspection, yet dense networks can collapse into ``hairballs'' in which connectivity is visible while scientifically meaningful structure remains difficult to discern. Analysts must determine which relationships matter, move between local detail and broader context, compare competing patterns, and connect visual observations to biological hypotheses. Supporting this work requires decisions about what to represent, aggregate, prioritize, and inspect next. How these decisions are made---and how they are distributed between scientists and computational systems---shapes both the course of exploration and the evidence available for scientific interpretation.

Computational assistance participates in these decisions by transforming representations, directing attention, and supporting the selection and execution of analytic actions. Clustering, aggregation, layout, and dimensionality reduction can make complex relationships tractable for inspection. Adaptive mechanisms can use interaction history or task context to recommend actions or reorganize what is shown. Conversational interfaces allow scientists to specify requests in language and, in more elaborate workflows, delegate sequences of operations for subsequent inspection. These mechanisms involve different forms of computational participation, rather than successive stages in which newer techniques necessarily provide better assistance. Established algorithms and newer AI methods are relevant here because of the work they perform within analysis. Their interventions shape what scientists encounter, which actions become available, and which decisions remain subject to direct human judgment.

The same analytic task can therefore involve markedly different divisions of work between scientist and system. To compare two biological subnetworks, an analyst might select both manually, inspect candidates ranked by an algorithm, respond to an adaptive recommendation, or ask a conversational system to identify and explain differences. The intent remains comparison, while responsibility for selecting candidates, executing operations, and interpreting their significance changes. Existing task frameworks provide essential vocabulary for describing such intent and its realization~\cite{Lee2006TaskTaxonomy,Brehmer2013WhyHowWhat}. Building on that vocabulary, understanding computational assistance also requires examining who initiates an operation, how its scope is determined, and how its results are assessed.

Examining this distribution of responsibility requires bringing together several perspectives in the review literature. Surveys of biological network visualization address the representations and interactions through which biological relationships become accessible~\cite{Ehlers2025}. Broader network-visualization surveys connect techniques with analytical tasks~\cite{Filipov2023Roadmap}, while reviews of immersive analytics and immersive network analysis examine visualization in spatial interaction environments~\cite{Fonnet2019ImmersiveAnalytics,Joos2025VisNetIA}. Our review brings these perspectives into a common comparison: how computational assistance participates in life-science analytic tasks across visualization environments, and what that participation implies for analyst control, interpretation, and validation. The focus is the relationship among task, computational intervention, and interaction environment, including situations in which several forms of assistance operate together.

We use \emph{locus of analytic agency} to describe how initiative and control over analytic actions, transformations, and judgments are distributed between scientists and computational systems. This distribution can change within a workflow: a scientist may delegate candidate selection, inspect the resulting representation, revise the selection criteria, and retain authority over biological interpretation and validation. Computational initiative does not itself establish the validity of a scientific claim. The lens therefore distinguishes the execution of an operation from decisions about its purpose, the interpretation of its output, and responsibility for accepting its consequences. We use agency as an interpretive lens for synthesis, rather than as an additional categorical assignment or a single score of automation.

The review examines this distribution through three complementary perspectives. \emph{Tasks} describe what analysts seek to accomplish: Navigation and Multiscale Orientation; Selection, Filtering, and Precision Interaction; Comparison and Differentiation; Sensemaking and Hypothesis Development; and Coordination and Collaborative Reasoning. \emph{Assistance} describes how computation participates, through four overlapping modes: Algorithmic, Adaptive, Conversational, and Immersive. \emph{Modality} describes the visualization and interaction environment: Desktop/Planar, Large Display, Virtual Reality (VR), Augmented/Mixed Reality (AR/MR), and CAVE. Tasks and assistance modes permit multiple labels; the categories do not partition systems into mutually exclusive classes. Immersive assistance specifically requires a computational mechanism that exploits spatial, embodied, or multisensory interaction in support of analysis. Merely displaying data in VR does not establish that mechanism.

Modality matters because the environment conditions how information is perceived, actions are expressed, and control is shared. Selecting a structure with a pointer, directing attention across a large display, and referring to a spatial object through speech and gesture place different demands on assistance. Spatial or embodied cues may also become inputs that a system interprets, introducing questions about ambiguous intent and correction. In collaborative settings, an intervention can affect both the requesting analyst and colleagues who share the representation. Modality thus constrains and enables the expression of agency without determining its distribution: the same environment may support direct manipulation, recommendations, delegated operations, or combinations of these relationships.

This STAR focuses on interactive visual analysis of dense relationships, derived structures, and multiscale biological context. It examines assistance within scientific workflows, including the transformations that make relationships inspectable and the mechanisms that help scientists steer exploration. Its scope does not extend to all computational biology or to predictive models without a relevant interactive visual-analytic role. The review is organized to connect mechanisms with the tasks they support, the human judgments they preserve or influence, and the evidence used to assess their consequences. A common search, screening, and multilabel coding procedure provides the methodological basis for that comparison, with human responsibility for decisions proposed through LLM-assisted processing.

\paragraph*{Research questions.}
We address four questions:
\begin{enumerate}
\item[\textbf{RQ1.}] What recurring analytic tasks characterize computationally assisted visual analysis in the life sciences?
\item[\textbf{RQ2.}] How do different forms of computational assistance participate in those tasks and redistribute analytic agency between human and system?
\item[\textbf{RQ3.}] How do visualization modalities condition the form, effectiveness, and consequences of that assistance?
\item[\textbf{RQ4.}] How are these systems evaluated, and what benefits, limitations, risks, and design obligations emerge as computational participation becomes more consequential?
\end{enumerate}

\paragraph*{Contributions.}
This STAR contributes (1) a reproducible systematic review of computationally assisted visual analytics for relational and multiscale life-science analysis; (2) a cross-cutting synthesis connecting analytic tasks, computational assistance, and visualization modalities; (3) an interpretation of these relationships through the locus of analytic agency; and (4) a comparative synthesis of evaluation evidence and research challenges.

Section~\ref{sec:methodology} describes the review methodology. Section~\ref{sec:foundations} establishes the conceptual foundations. Sections~\ref{sec:tasks}--\ref{sec:modality} develop the task, assistance, and modality syntheses; Section~\ref{sec:evaluation} examines evaluation evidence. Section~\ref{sec:challenges} connects the resulting tensions to research opportunities, and Section~\ref{sec:conclusion} concludes.

\draftnote{Complete prose draft; methodological and contribution statements describe the intended completed report. Confirm them against the executed review before submission. Citations are retained from the supplied bibliography and require bibliographic/source checking during integration. Foundational references are now discussed in Sections 3.5--3.6. Complete the prior-survey comparison table before finalizing the scope of the survey-gap claim. No corpus prevalence, final counts, or completed validation results are asserted.}
